\section{Trabalhos Relacionados}

A detecção de anomalias com autoencoders recorrentes baseia-se na ideia de que a
reconstrução de padrões aprendidos falha sobre observações fora da distribuição de
treino. \citet{li2020lstm}, em referência amplamente reproduzida, aplica um
LSTM-Autoencoder a preços de ações (Johnson \& Johnson), adotando janelas de 30 passos
de tempo e definindo o limiar de anomalia a partir do erro de reconstrução do treino.
\citet{valkov2019lstm} fornece uma implementação canônica em Keras, com escolhas de
projeto que servem de referência prática (camadas LSTM de 64 unidades, \emph{dropout}
de 0{,}2, \texttt{validation\_split} de 0{,}1 e, crucialmente, \texttt{shuffle=False}
para preservar a ordem temporal). \citet{petrovic2019lstm} disponibiliza um repositório
com abordagem equivalente.

Trabalhos mais recentes ampliam o desenho experimental. \citet{liu2025robust} propõem um
arcabouço híbrido que acopla um LSTM-Autoencoder a redes adversárias generativas e
introduzem explicitamente um \emph{mecanismo de injeção artificial de anomalias} para
avaliação quantitativa, validado em múltiplas categorias de ações --- desenho que
inspira diretamente a avaliação setorial e por injeção adotada neste projeto. No domínio
de criptomoedas, \citet{bitcoin2021anomaly} exploram a mesma família de modelos sobre
valores do Bitcoin.

Em relação a esses trabalhos, a principal lacuna endereçada aqui é a aplicação ao mercado
brasileiro (B3) com correlação documentada a eventos nacionais e comparação explícita
entre estratégias de limiar estático e dinâmico.

% \citet usado acima requer natbib; caso o estilo não suporte, ver nota em report/README.
